%==============================================================
%  estilos/listas.tex
%  Configuración de listas con enumitem, coherente con la convención de
%  párrafos de la tesis: sangría en todos los párrafos y ningún espacio
%  vertical entre ellos (ver «Formato de párrafo» en tesis.tex).
%
%  Se carga desde el preámbulo con  %==============================================================
%  estilos/listas.tex
%  Configuración de listas con enumitem, coherente con la convención de
%  párrafos de la tesis: sangría en todos los párrafos y ningún espacio
%  vertical entre ellos (ver «Formato de párrafo» en tesis.tex).
%
%  Se carga desde el preámbulo con  %==============================================================
%  estilos/listas.tex
%  Configuración de listas con enumitem, coherente con la convención de
%  párrafos de la tesis: sangría en todos los párrafos y ningún espacio
%  vertical entre ellos (ver «Formato de párrafo» en tesis.tex).
%
%  Se carga desde el preámbulo con  %==============================================================
%  estilos/listas.tex
%  Configuración de listas con enumitem, coherente con la convención de
%  párrafos de la tesis: sangría en todos los párrafos y ningún espacio
%  vertical entre ellos (ver «Formato de párrafo» en tesis.tex).
%
%  Se carga desde el preámbulo con  \input{estilos/listas}
%
%  PARA VOLVER AL ASPECTO POR DEFECTO DE LaTeX: comenta esa línea en
%  tesis.tex. Ninguna otra parte del documento depende de este archivo.
%
%  Medido en el PDF con la geometría de la tesis (pdftotext -bbox,
%  márgenes de 96.38 pt y sangría de párrafo de 112.65 pt):
%    antes y después del bloque:  11.95 pt por defecto  ->  6.77 pt
%    entre reactivos:              8.96 pt por defecto  ->  3.39 pt
%    cuerpo del reactivo:        123.65 pt por defecto  ->  112.74 pt
%                                (= la sangría del párrafo, 1.5em)
%==============================================================

%--------------------------------------------------------------
%  ESPACIADO VERTICAL
%  parsep = 0 porque en este documento no hay espacio entre párrafos;
%  media línea antes y después del bloque para separarlo del texto, y un
%  cuarto de línea entre reactivos. partopsep = 0 evita el espacio doble
%  cuando la lista abre un párrafo nuevo.
%--------------------------------------------------------------
\setlist{
  topsep    = 0.5\baselineskip,
  partopsep = 0pt,
  itemsep   = 0.25\baselineskip,
  parsep    = 0pt,
}

%--------------------------------------------------------------
%  INDENTACIÓN HORIZONTAL (itemize y enumerate)
%  El cuerpo del reactivo se alinea con la sangría de los párrafos y el
%  marcador cuelga en el margen. Está atado a \parindent (1.5em), así que
%  si algún día cambia la sangría, la lista la sigue sola.
%--------------------------------------------------------------
\setlist[itemize,enumerate]{
  labelindent = 0pt,
  leftmargin  = \parindent,
  labelsep    = 0.5em,
}

% Sublistas más compactas: menos aire arriba y entre reactivos.
\setlist[itemize,enumerate,2]{
  topsep     = 0.25\baselineskip,
  itemsep    = 0.1\baselineskip,
  leftmargin = *,
}

% Las descripciones heredan el espaciado vertical de arriba y conservan su
% formato por defecto (rótulo en negrita seguido del texto). Si algún día
% se quiere el cuerpo alineado como en itemize y enumerate, añade aquí:
%   \setlist[description]{leftmargin=\parindent, labelsep=0.5em}

%--------------------------------------------------------------
%  LISTAS EN LÍNEA (enumerate*, itemize*)
%  Requieren que enumitem se cargue con la opción `inline`, que ya está en
%  el preámbulo. Separadores en español: «(i) A, (ii) B y (iii) C».
%
%  Dos trampas de sintaxis, ambas comprobadas:
%   - La configuración va con ASTERISCO en \setlist, no en el nombre del
%     entorno:  \setlist*[enumerate]{...}   (no \setlist[enumerate*]{...},
%     que da «Missing number, treated as zero»).
%   - El espacio inicial de un valor se pierde si va entre llaves, así que
%     « y » hay que escribirlo «\ y\ » (con espacio de control).
%
%  Si los reactivos llevan comas internas, conviene separarlos con punto y
%  coma; pásalo localmente en esa lista:
%      \begin{enumerate*}[label=(\roman*), itemjoin={;\ }, itemjoin*={;\ y\ }]
%--------------------------------------------------------------
\setlist*[enumerate]{itemjoin={,\ }, itemjoin*={\ y\ }}
\setlist*[itemize]{itemjoin={,\ }, itemjoin*={\ y\ }}

%
%  PARA VOLVER AL ASPECTO POR DEFECTO DE LaTeX: comenta esa línea en
%  tesis.tex. Ninguna otra parte del documento depende de este archivo.
%
%  Medido en el PDF con la geometría de la tesis (pdftotext -bbox,
%  márgenes de 96.38 pt y sangría de párrafo de 112.65 pt):
%    antes y después del bloque:  11.95 pt por defecto  ->  6.77 pt
%    entre reactivos:              8.96 pt por defecto  ->  3.39 pt
%    cuerpo del reactivo:        123.65 pt por defecto  ->  112.74 pt
%                                (= la sangría del párrafo, 1.5em)
%==============================================================

%--------------------------------------------------------------
%  ESPACIADO VERTICAL
%  parsep = 0 porque en este documento no hay espacio entre párrafos;
%  media línea antes y después del bloque para separarlo del texto, y un
%  cuarto de línea entre reactivos. partopsep = 0 evita el espacio doble
%  cuando la lista abre un párrafo nuevo.
%--------------------------------------------------------------
\setlist{
  topsep    = 0.5\baselineskip,
  partopsep = 0pt,
  itemsep   = 0.25\baselineskip,
  parsep    = 0pt,
}

%--------------------------------------------------------------
%  INDENTACIÓN HORIZONTAL (itemize y enumerate)
%  El cuerpo del reactivo se alinea con la sangría de los párrafos y el
%  marcador cuelga en el margen. Está atado a \parindent (1.5em), así que
%  si algún día cambia la sangría, la lista la sigue sola.
%--------------------------------------------------------------
\setlist[itemize,enumerate]{
  labelindent = 0pt,
  leftmargin  = \parindent,
  labelsep    = 0.5em,
}

% Sublistas más compactas: menos aire arriba y entre reactivos.
\setlist[itemize,enumerate,2]{
  topsep     = 0.25\baselineskip,
  itemsep    = 0.1\baselineskip,
  leftmargin = *,
}

% Las descripciones heredan el espaciado vertical de arriba y conservan su
% formato por defecto (rótulo en negrita seguido del texto). Si algún día
% se quiere el cuerpo alineado como en itemize y enumerate, añade aquí:
%   \setlist[description]{leftmargin=\parindent, labelsep=0.5em}

%--------------------------------------------------------------
%  LISTAS EN LÍNEA (enumerate*, itemize*)
%  Requieren que enumitem se cargue con la opción `inline`, que ya está en
%  el preámbulo. Separadores en español: «(i) A, (ii) B y (iii) C».
%
%  Dos trampas de sintaxis, ambas comprobadas:
%   - La configuración va con ASTERISCO en \setlist, no en el nombre del
%     entorno:  \setlist*[enumerate]{...}   (no \setlist[enumerate*]{...},
%     que da «Missing number, treated as zero»).
%   - El espacio inicial de un valor se pierde si va entre llaves, así que
%     « y » hay que escribirlo «\ y\ » (con espacio de control).
%
%  Si los reactivos llevan comas internas, conviene separarlos con punto y
%  coma; pásalo localmente en esa lista:
%      \begin{enumerate*}[label=(\roman*), itemjoin={;\ }, itemjoin*={;\ y\ }]
%--------------------------------------------------------------
\setlist*[enumerate]{itemjoin={,\ }, itemjoin*={\ y\ }}
\setlist*[itemize]{itemjoin={,\ }, itemjoin*={\ y\ }}

%
%  PARA VOLVER AL ASPECTO POR DEFECTO DE LaTeX: comenta esa línea en
%  tesis.tex. Ninguna otra parte del documento depende de este archivo.
%
%  Medido en el PDF con la geometría de la tesis (pdftotext -bbox,
%  márgenes de 96.38 pt y sangría de párrafo de 112.65 pt):
%    antes y después del bloque:  11.95 pt por defecto  ->  6.77 pt
%    entre reactivos:              8.96 pt por defecto  ->  3.39 pt
%    cuerpo del reactivo:        123.65 pt por defecto  ->  112.74 pt
%                                (= la sangría del párrafo, 1.5em)
%==============================================================

%--------------------------------------------------------------
%  ESPACIADO VERTICAL
%  parsep = 0 porque en este documento no hay espacio entre párrafos;
%  media línea antes y después del bloque para separarlo del texto, y un
%  cuarto de línea entre reactivos. partopsep = 0 evita el espacio doble
%  cuando la lista abre un párrafo nuevo.
%--------------------------------------------------------------
\setlist{
  topsep    = 0.5\baselineskip,
  partopsep = 0pt,
  itemsep   = 0.25\baselineskip,
  parsep    = 0pt,
}

%--------------------------------------------------------------
%  INDENTACIÓN HORIZONTAL (itemize y enumerate)
%  El cuerpo del reactivo se alinea con la sangría de los párrafos y el
%  marcador cuelga en el margen. Está atado a \parindent (1.5em), así que
%  si algún día cambia la sangría, la lista la sigue sola.
%--------------------------------------------------------------
\setlist[itemize,enumerate]{
  labelindent = 0pt,
  leftmargin  = \parindent,
  labelsep    = 0.5em,
}

% Sublistas más compactas: menos aire arriba y entre reactivos.
\setlist[itemize,enumerate,2]{
  topsep     = 0.25\baselineskip,
  itemsep    = 0.1\baselineskip,
  leftmargin = *,
}

% Las descripciones heredan el espaciado vertical de arriba y conservan su
% formato por defecto (rótulo en negrita seguido del texto). Si algún día
% se quiere el cuerpo alineado como en itemize y enumerate, añade aquí:
%   \setlist[description]{leftmargin=\parindent, labelsep=0.5em}

%--------------------------------------------------------------
%  LISTAS EN LÍNEA (enumerate*, itemize*)
%  Requieren que enumitem se cargue con la opción `inline`, que ya está en
%  el preámbulo. Separadores en español: «(i) A, (ii) B y (iii) C».
%
%  Dos trampas de sintaxis, ambas comprobadas:
%   - La configuración va con ASTERISCO en \setlist, no en el nombre del
%     entorno:  \setlist*[enumerate]{...}   (no \setlist[enumerate*]{...},
%     que da «Missing number, treated as zero»).
%   - El espacio inicial de un valor se pierde si va entre llaves, así que
%     « y » hay que escribirlo «\ y\ » (con espacio de control).
%
%  Si los reactivos llevan comas internas, conviene separarlos con punto y
%  coma; pásalo localmente en esa lista:
%      \begin{enumerate*}[label=(\roman*), itemjoin={;\ }, itemjoin*={;\ y\ }]
%--------------------------------------------------------------
\setlist*[enumerate]{itemjoin={,\ }, itemjoin*={\ y\ }}
\setlist*[itemize]{itemjoin={,\ }, itemjoin*={\ y\ }}

%
%  PARA VOLVER AL ASPECTO POR DEFECTO DE LaTeX: comenta esa línea en
%  tesis.tex. Ninguna otra parte del documento depende de este archivo.
%
%  Medido en el PDF con la geometría de la tesis (pdftotext -bbox,
%  márgenes de 96.38 pt y sangría de párrafo de 112.65 pt):
%    antes y después del bloque:  11.95 pt por defecto  ->  6.77 pt
%    entre reactivos:              8.96 pt por defecto  ->  3.39 pt
%    cuerpo del reactivo:        123.65 pt por defecto  ->  112.74 pt
%                                (= la sangría del párrafo, 1.5em)
%==============================================================

%--------------------------------------------------------------
%  ESPACIADO VERTICAL
%  parsep = 0 porque en este documento no hay espacio entre párrafos;
%  media línea antes y después del bloque para separarlo del texto, y un
%  cuarto de línea entre reactivos. partopsep = 0 evita el espacio doble
%  cuando la lista abre un párrafo nuevo.
%--------------------------------------------------------------
\setlist{
  topsep    = 0.5\baselineskip,
  partopsep = 0pt,
  itemsep   = 0.25\baselineskip,
  parsep    = 0pt,
}

%--------------------------------------------------------------
%  INDENTACIÓN HORIZONTAL (itemize y enumerate)
%  El cuerpo del reactivo se alinea con la sangría de los párrafos y el
%  marcador cuelga en el margen. Está atado a \parindent (1.5em), así que
%  si algún día cambia la sangría, la lista la sigue sola.
%--------------------------------------------------------------
\setlist[itemize,enumerate]{
  labelindent = 0pt,
  leftmargin  = \parindent,
  labelsep    = 0.5em,
}

% Sublistas más compactas: menos aire arriba y entre reactivos.
\setlist[itemize,enumerate,2]{
  topsep     = 0.25\baselineskip,
  itemsep    = 0.1\baselineskip,
  leftmargin = *,
}

% Las descripciones heredan el espaciado vertical de arriba y conservan su
% formato por defecto (rótulo en negrita seguido del texto). Si algún día
% se quiere el cuerpo alineado como en itemize y enumerate, añade aquí:
%   \setlist[description]{leftmargin=\parindent, labelsep=0.5em}

%--------------------------------------------------------------
%  LISTAS EN LÍNEA (enumerate*, itemize*)
%  Requieren que enumitem se cargue con la opción `inline`, que ya está en
%  el preámbulo. Separadores en español: «(i) A, (ii) B y (iii) C».
%
%  Dos trampas de sintaxis, ambas comprobadas:
%   - La configuración va con ASTERISCO en \setlist, no en el nombre del
%     entorno:  \setlist*[enumerate]{...}   (no \setlist[enumerate*]{...},
%     que da «Missing number, treated as zero»).
%   - El espacio inicial de un valor se pierde si va entre llaves, así que
%     « y » hay que escribirlo «\ y\ » (con espacio de control).
%
%  Si los reactivos llevan comas internas, conviene separarlos con punto y
%  coma; pásalo localmente en esa lista:
%      \begin{enumerate*}[label=(\roman*), itemjoin={;\ }, itemjoin*={;\ y\ }]
%--------------------------------------------------------------
\setlist*[enumerate]{itemjoin={,\ }, itemjoin*={\ y\ }}
\setlist*[itemize]{itemjoin={,\ }, itemjoin*={\ y\ }}
